\section{相关技术与需求分析}

\subsection{相关技术基础}

\subsubsection{大语言模型与结构化需求分析}

大语言模型擅长从非结构化文本中抽取实体、关系和约束条件。在本项目中，模型不直接承担 CAD 脚本生成，而是只承担“把自然语言需求转成标准 JSON Plan”的任务。为了保证输出具有可消费性，系统在提示词中显式限定了 JSON 格式、允许的零件形状、参数字段名称以及关系类型。

\subsubsection{参数化 CAD 与 FeatureScript}

参数化 CAD 的核心在于通过参数、特征和几何关系描述模型，而非只保留静态几何结果。Onshape 的 FeatureScript 机制支持以代码形式定义特征，包括立方体、圆柱体、布尔运算和属性设置等操作 \cite{onshape_fsdoc}。与传统 GUI 建模相比，脚本化建模更适合进行自动生成、模板化复用和版本管理。因此，本文将经过校验的装配 Plan 翻译为 FeatureScript，作为系统输出结果。

\subsubsection{强类型数据校验}

由于自然语言分析存在不确定性，系统必须在生成阶段之前对中间结果进行严格校验。项目使用 Pydantic 风格的强类型模型对装配名、零件 ID、形状类别、参数集合、关系类型和数值有效性进行检查。例如，\texttt{flange} 形状要求法兰直径必须大于轴直径，\texttt{tapered\_cylinder} 必须同时给出上下直径和高度。通过在中间层进行约束收缩，可以避免错误继续扩散到最终脚本。

\subsection{业务场景分析}

本系统面向的不是复杂工业装配，而是结构较简单、逻辑较清晰、参数相对完整的教学与演示场景，例如模具主件、简单夹具、法兰支撑组件和导向柱套类结构等。其典型输入并非标准 CAD 操作指令，而是自然语言描述，如“设计一个冷拉延模具，只建五个主件，上下模座尺寸为多少，凸模外径与高度为多少”等。

根据仓库中的示例需求文本，用户输入通常包含工件或装配对象名称、关键零件及其功能、主要几何尺寸、总体尺寸约束，以及是否省略某些复杂构件等信息。这类需求具备较强语义性，但并不直接等价于 CAD 操作序列，因此系统需要先将语义需求转化为统一的工程中间表示，再进入脚本生成阶段。

\subsection{功能与非功能需求分析}

结合项目目标与当前实现，系统需满足以下要求：能够读取文本输入并抽取零件实体、形状类别、尺寸参数和装配关系，输出统一格式的 JSON Plan；能够对 Plan 进行严格校验，尽早拦截非法字段、错误命名和不完整参数集合；能够基于合法 Plan 生成对应的 FeatureScript 文件，并在局部失败时保留错误注释；能够同时支持 CLI 与本地浏览器界面，便于脚本化使用和答辩演示。

在非功能方面，系统需要兼顾可控性、可验证性、可扩展性和易用性。具体而言，输出必须可预测且可追踪，关键模块应具有测试支撑，系统结构应支持扩展更多零件形状和输出模板，交互方式则应保持低学习成本。
